\section{Proper scoring rule 与校准衡量概率是否诚实}
\label{sec:11-Mathematical-Statistics/07-Statistical-Decision-and-Reliable-Prediction/04}

按上一节算出的阈值 \(0.038\) 部署两周，送检量是测算值的三倍，而送检者里的阳性比例只有报告概率所暗示的三分之一。模型没坏，AUC 与离线一致。翻查后发现：被模型标成 \(0.04\) 的那批人，实际发生率约 \(0.013\)。

阈值公式要求 \(p\) 是校准过的概率，这一节检查这个前提。要问两件事：模型有没有动机把自己的信念如实报出来，以及报出来的数字在长期频率上兑不兑现。

\subsection{只看排序指标，模型没有说真话的动机}

设 \(Y\in\{0,1\}\)，真实条件概率 \(p(x)=P(Y=1\given X=x)\)，模型报告 \(q(x)\in[0,1]\)。评分规则 \(S(q,Y)\) 给这份报告与实际结果配一个损失。

AUC 对分数的任何单调变换都不变。把所有 logit 乘以 \(3\)，排序一模一样，AUC 一个小数点都不动，概率却全被推向 \(0\) 与 \(1\)。因此以 AUC 为唯一验收指标时，概率的数值是自由的，模型没有理由报准。要约束数值，得换一类评分。

Brier 损失是最简单的一个：

\[
S_{\mathrm{Brier}}(q,Y)=(Y-q)^2 .
\]

给定 \(X=x\)，它的条件期望可以拆开：

\[
\E\!\left[(Y-q)^2\given X=x\right]=p(x)\bigl(1-p(x)\bigr)+\bigl(p(x)-q\bigr)^2 .
\]

第一项是知道 \(x\) 之后仍无法消除的 Bernoulli 波动，与报告无关；第二项在 \(q=p(x)\) 时取零。于是诚实报告是唯一的期望最优策略。凡是具备这个性质的评分规则称为严格 proper，它把``说真话''变成模型自身的最优解，而不是一条附加要求。

对数损失

\[
S_{\log}(q,Y)=-Y\log q-(1-Y)\log(1-q)
\]

同样严格 proper，它的超额风险恰是真实 Bernoulli 分布到报告分布的 KL 散度，这一点在\hyperref[sec:07-Information-Theory/03-Divergences/01]{``错误模型的代价：从交叉熵到 KL 散度''}里已经从编码长度的角度算过一遍。两者的差别在尾部：对数损失对``报 \(0.001\) 却发生了''施加无界惩罚，Brier 有界，对极端错误宽容得多。安全关键的场景倾向前者，评分被少数离群样本主导的场景倾向后者。

\subsection{校准是同一预测值上的长期频率对得上}

记随机预测 \(Q=q(X)\)。完美校准要求

\[
\E[Y\given Q]=Q
\]

几乎处处成立：所有报 \(0.7\) 的对象中，长期约有 \(70\%\) 发生。检查它的标准工具是可靠性图。

\begin{center}
\begin{tikzpicture}[font=\footnotesize,>=stealth,x=4.2cm,y=3.2cm]
  \draw[black!30] (0,0) rectangle (1,1);
  \draw[black!50,dashed] (0,0) -- (1,1);
  \node[black!55,rotate=37,anchor=south] at (0.66,0.66) {完美校准};
  \draw[black!85,thick] plot[smooth] coordinates
    {(0.05,0.02) (0.15,0.07) (0.25,0.13) (0.35,0.21) (0.45,0.30)
     (0.55,0.39) (0.65,0.50) (0.75,0.60) (0.85,0.72) (0.95,0.86)};
  \foreach \a/\b in {0.05/0.02,0.15/0.07,0.25/0.13,0.35/0.21,0.45/0.30,
                     0.55/0.39,0.65/0.50,0.75/0.60,0.85/0.72,0.95/0.86}
    \node[circle,fill=black!85,inner sep=1.1pt] at (\a,\b) {};
  \draw[->,black!70] (-0.02,0) -- (1.14,0);
  \draw[->,black!70] (0,-0.02) -- (0,1.12) node[left] {实际频率};
  \node[black!70,anchor=north] at (0.45,-0.32) {预测概率 \(q\)};
  \foreach \a/\h in {0.05/0.17,0.15/0.13,0.25/0.10,0.35/0.075,0.45/0.06,
                     0.55/0.05,0.65/0.04,0.75/0.032,0.85/0.024,0.95/0.016}
    \draw[black!35,line width=3.2pt] (\a,-0.07) -- (\a,{-0.07-\h});
  \node[black!55,anchor=west] at (1.04,-0.16) {每箱样本量};
\end{tikzpicture}

\emph{可靠性图：十个箱内的预测概率对实际频率。点整体落在对角线下方，模型系统性报高；下方细柱是每箱样本量，右侧几箱太少，那里的偏离读不出可信的结论。}
\end{center}

图右上角那几个点偏离最大，也最不该被当真——柱子只有左端的十分之一，几十个样本上的频率本来就抖。这提示一套可以照做的画法：用没参与训练也没参与调阈值的独立数据；分 \(10\) 到 \(20\) 个箱，样本不多时改用等频分箱以免右侧箱空；每个点配一条 Wilson 区间的竖线；把每箱样本量画在下方或直接标出来；关键子群各画一张。只报一个期望校准误差数字则要小心，它随分箱边界、箱数和加权方式变化，而且在同一份数据上先重校准再评估必然乐观。

校准声明还捆着三样东西：预测时点、目标人群、后续政策。急诊模型说某组患者二十四小时恶化概率 \(0.2\)，要求的是在同样时点、同样人群、同样分诊政策的大量重复中约五分之一恶化。它不要求某天的五位患者里恰好一位恶化，也不允许从``全院总体校准''推出``夜班患者校准''——换一种分组就换了条件事件，声明的对象跟着变。更麻烦的是，如果这个概率会触发提前干预，之后观测到的恶化率已经被干预压低了，此时必须讲清校准针对的是``当前政策下的结果''还是``不干预时的自然风险''，两者的生成机制不同，相关讨论见\hyperref[sec:13-Machine-Learning/12-Causal-Inference/01]{``观察、干预与反事实回答不同问题''}。

\subsection{区分度与校准互不蕴含}

一个永远输出总体发生率 \(P(Y=1)\) 的常数模型完美校准，它对谁都说同一个数，没有任何用处。反过来，前面那个 logit 乘 \(3\) 的模型排序完好而概率全乱。两种性质可以单独成立，也可以单独失败。

Brier 分数把它们摆在同一个式子里。把预测分成 \(K\) 组，第 \(k\) 组样本量 \(n_k\)、平均预测 \(\bar q_k\)、实际频率 \(\bar y_k\)，总体频率 \(\bar y\)，则

\[
\frac1n\sum_{i=1}^n (y_i-q_i)^2
=\underbrace{\bar y(1-\bar y)}_{\text{不确定性}}
-\underbrace{\frac1n\sum_{k}n_k(\bar y_k-\bar y)^2}_{\text{区分度}}
+\underbrace{\frac1n\sum_{k}n_k(\bar q_k-\bar y_k)^2}_{\text{可靠性}} .
\]

第一项只与任务本身有关，模型碰不到。区分度带负号，衡量各组的实际发生率彼此拉开了多远——这是模型提供的信息量。可靠性带正号，衡量各组的报告值离本组实际频率有多远，正是可靠性图里点到对角线的竖直距离的加权平方和。一个分数变好，可能来自信息变多，也可能来自数值修准，还可能来自换了个更容易的测试总体，所以跨数据集直接比 Brier 或对数损失要连基准率一起看。把三项分开报，改进来自哪里就藏不住了。

\subsection{重校准修数值，不补信息}

Platt scaling、isotonic regression、temperature scaling 都是在独立校准集上学一个从原分数到概率的映射。它们能压平系统性的过度自信，也能把可靠性项打下来，但区分度项那一部分它们动不了：单调映射不改变排序，模型从未捕捉到的区别不会因为重校准而出现。分布一变，学到的映射也要重学。

有几种常见训练手法会让重校准变成必需项。类别加权、过采样、focal loss、label smoothing 都改变了训练目标对应的分布，最优输出不再等于自然部署人群里的条件概率。做法是先说清训练损失对应哪个目标分布，再在真实目标分布的独立样本上单独校准并验证，两份数据不能是同一份。

值得记住的是，前面那套阈值公式和这一节的校准检查配套使用才成立。阈值 \(t^\ast\) 的推导默认 \(p\) 的字面含义可信；概率整体报高时，同样的 \(t^\ast\) 会放进来一大批实际风险很低的对象，就是开头那个送检量三倍的现象。修法有两条并且互不冲突：把概率重校准，或者在原分数上直接按目标送检量反解阈值，前者让数字恢复含义，后者绕开数字只保排序。

\subsection{诚实的概率仍不是可交付的保证}

即使校准完美，模型交给下游的仍是一个点概率，它回答的是``这一个对象的发生率是多少''，而不是``这一批结论里有多少会落空''。有些场合需要后一种承诺：给出一个集合或区间，并说明它以多大比例盖住真值，且这句话在当前样本量下就成立，不靠渐近。下一节用一个只需要可交换性的构造给出这种保证。
